\section{Experimental Results}

The evaluation of our multimodal fusion framework was conducted on a 20\% hold-out test set derived from the 500 000 synthetic patient records. This section details the validation of the synthetic data, the predictive performance, the computational efficiency, and a deep dive into model interpretability.

\subsection{Validation of Synthetic Physiological Coherence}
Before evaluating the predictive models, it is crucial to verify that the synthetic generation engine produced realistic clinical relationships rather than random statistical noise. Figure \ref{fig:corr_matrix} presents the physiological correlation matrix of the generated variables.

\begin{figure}[htbp]
\centerline{\includegraphics[width=0.48\textwidth]{10_Correlation_Matrix.png}}
\caption{Physiological Correlation Matrix confirming the realistic inter-variable dependencies within the synthetic IoMT dataset.}
\label{fig:corr_matrix}
\end{figure}

The heatmap demonstrates strong positive and negative correlation clusters that align with medical reality. For instance, the correlation block between distinct biomarkers (such as Leukocytes\_G\_L and CRP\_mg\_L) confirms that our rule-based system successfully maintained known physiological dependencies, rendering the dataset a valid, highly structured benchmark for medical Deep Learning.

\subsection{Experimental Setup and Hyperparameters}
To ensure reproducibility, the hyperparameter configurations for both architectures are summarized in Table \ref{tab:hyperparams}. Both models were implemented in PyTorch and trained on an optimized hardware environment.

\begin{table}[htbp]
\caption{Model Training Hyperparameters}
\begin{center}
\begin{tabular}{|l|c|c|}
\hline
\textbf{Hyperparameter} & \textbf{MLP Baseline} & \textbf{Tabular Transformer} \\
\hline
Optimizer & AdamW & AdamW \\
Learning Rate & 0.002 & 0.002 \\
Weight Decay & $1e-4$ & $1e-4$ \\
Dropout Rate & 0.3 / 0.2 & 0.2 \\
Batch Size & 1024 & 1024 \\
Embedding Dim ($d_{model}$) & N/A & 64 \\
Attention Heads & N/A & 4 \\
\hline
\end{tabular}
\label{tab:hyperparams}
\end{center}
\end{table}

\subsection{Predictive Performance and ROC-AUC Analysis}
As illustrated in Figure \ref{fig:perf_global}, both models demonstrate high efficiency in classification tasks across the 24 highly imbalanced pathologies.

\begin{figure}[htbp]
\centerline{\includegraphics[width=0.48\textwidth]{01_Global_Performances.png}}
\caption{Global Performance metrics illustrating the trade-offs between Accuracy and Precision across models.}
\label{fig:perf_global}
\end{figure}

The Tabular Transformer achieved an impressive Global Accuracy of 77.1\%, outperforming the Neural Network baseline (75.2\%). However, a deeper look at the metrics reveals an interesting dynamic: the Neural Network achieved a significantly better Precision (66.5\% vs 59.1\%) and a higher Macro F1-Score (63.5\% vs 60.8\%). This suggests that while the Transformer is highly effective at categorizing the majority classes (boosting overall accuracy), the simpler MLP, aided by class weighting, was slightly more conservative and precise when predicting minority target pathologies.

The discriminative power is further validated by the Receiver Operating Characteristic (ROC) analysis in Figure \ref{fig:roc}.

\begin{figure}[htbp]
\centerline{\includegraphics[width=0.48\textwidth]{05_ROC_Curves.png}}
\caption{Comparative ROC Curves demonstrating near-perfect sensitivity and specificity for both models (AUC = 0.994).}
\label{fig:roc}
\end{figure}

Both architectures reached an identical, near-perfect Micro-Average AUC of 0.994. This extremely high score proves that the mathematical signatures generated by our synthetic engine are highly distinct. The models are not randomly guessing but have successfully mapped the multidimensional boundaries of the 24 diseases.

\subsection{Precision-Recall Analysis for Imbalanced Classes}
While ROC-AUC is standard, the Precision-Recall (PR) curve provides a more stringent assessment for heavily imbalanced medical datasets.

\begin{figure}[htbp]
\centerline{\includegraphics[width=0.48\textwidth]{09_Precision_Recall_Curve.png}}
\caption{Precision-Recall (PR) Curve illustrating model robustness on imbalanced clinical classes.}
\label{fig:pr_curve}
\end{figure}

Figure \ref{fig:pr_curve} demonstrates the PR curves for both architectures. Both the Tabular Transformer and the Neural Network achieved an identical Micro-Average Precision (AP) of 0.906. The curves remain remarkably high across the recall spectrum, confirming that both models can identify a vast majority of critical cases (high sensitivity) without triggering an excessive number of false alarms (high precision).

\subsection{Computational Overhead and Scalability}
A major finding of this study relates to the severe trade-off between architectural complexity and computational cost, as highlighted in Figure \ref{fig:temps}.

\begin{figure}[htbp]
\centerline{\includegraphics[width=0.48\textwidth]{02_Computation_Time.png}}
\caption{Total training time (left) and inference latency (right) comparison.}
\label{fig:temps}
\end{figure}

The Transformer required 164.0 seconds to train, compared to a mere 34.4 seconds for the MLP baseline. This discrepancy is mirrored in inference time: 0.695 seconds for the Transformer versus 0.302 seconds for the MLP. The $O(N^2)$ algorithmic complexity of the Multi-Head Attention mechanism introduces a massive computational penalty. Therefore, while Transformers are powerful for cloud-based hospital servers, the MLP remains the superior architecture for Edge computing and real-time IoMT deployment.

\subsection{Convergence and Overfitting Control}
The stability of the training process was monitored through the Cross-Entropy loss curves shown in Figure \ref{fig:loss}.

\begin{figure*}[t]
\centering
\includegraphics[width=0.8\textwidth]{07_Learning_Curves.png}
\caption{Learning curves for MLP (left) and Transformer (right) architectures showing stable loss reduction without overfitting.}
\label{fig:loss}
\end{figure*}

Both architectures demonstrate healthy convergence. The Neural Network (left) converged smoothly over approximately 24 epochs. The Transformer (right) learned aggressively, triggering the Early Stopping mechanism earlier (around 18 epochs). Most importantly, for both models, the Validation Loss (dashed red line) closely tracks or remains below the Training Loss (solid blue line). This confirms that our dropout strategies and Early Stopping successfully prevented the models from memorizing the synthetic dataset.

\subsection{Error Analysis and Clinical Overlap}
To understand the limitations of the fusion logic, we analyzed the confusion matrices (Figures \ref{fig:cm_nn} and \ref{fig:cm_tf}).

\begin{figure}[htbp]
\centerline{\includegraphics[width=0.48\textwidth]{03_Confusion_Matrix_NN.png}}
\caption{Confusion Matrix for the Neural Network baseline.}
\label{fig:cm_nn}
\end{figure}

\begin{figure}[htbp]
\centerline{\includegraphics[width=0.48\textwidth]{04_Confusion_Matrix_Transformer.png}}
\caption{Confusion Matrix for the Tabular Transformer.}
\label{fig:cm_tf}
\end{figure}

Both matrices show an exceptionally strong diagonal, confirming high overall accuracy. However, deliberate clinical noise programmed into our synthetic data induced logical misclassifications that reflect true medical complexity. For example, both models exhibit hesitation between overlapping respiratory and infectious conditions; notably, cases of Asthma\_Attack and Bacterial\_Pneumonia are frequently misclassified as Anxiety\_Hyperventilation due to shared vital sign anomalies. Additionally, a distinct diagnostic overlap emerges in the neurological branch, where both architectures show significant cross-confusion between Severe\_Migraine and Epileptic\_Seizure. This validates that the synthetic dataset is non-trivial and correctly mimics the physiological ambiguity and overlapping phenotypes found in real-world clinical environments.

\subsection{Model Interpretability and Feature Importance}
To ensure that our Deep Learning models did not merely memorize spurious statistical artifacts, we applied SHapley Additive exPlanations (SHAP) to interpret the decision-making process. Figure \ref{fig:shap} presents the feature importance and interaction values for a subset of the test data.

\begin{figure*}[t]
\centering
\includegraphics[width=0.8\textwidth]{06_SHAP_Explicability.png}
\caption{SHAP Summary Plot demonstrating the global feature importance and interaction values.}
\label{fig:shap}
\end{figure*}

The distinct clustering confirms that the model relies on true physiological biomarkers to drive its predictions. The visual distribution of SHAP values provides strong clinical validation of our synthetic engine. For instance, features such as Obs\_ChestPain, ECG\_Arrhythmia, and Hypertensive exhibit massive SHAP interaction values, proving that the neural network correctly identified them as the primary drivers for cardiological conditions. Similarly, the presence of physiological stress indicators like Obs\_Sweating, Obs\_Chills, and elevated Heart Rate (HR) strongly dictates the decision boundaries for other emergencies. Ultimately, this explainability analysis confirms that the predictive baselines successfully learned the expert-driven physiological rules embedded within the synthetic data generation pipeline, rather than relying on random noise.